\documentclass[a4paper,12pt]{article}
\usepackage[T2A]{fontenc}
\usepackage[utf8]{inputenc}
\usepackage[russian]{babel}
\usepackage{amsmath,amssymb,mathtools}
\usepackage{geometry}
\geometry{left=24mm,right=24mm,top=21mm,bottom=24mm}
\usepackage{microtype}
\usepackage{icomma}
\usepackage{xcolor}
\usepackage{tikz}
\usetikzlibrary{arrows.meta,positioning,shapes.geometric,calc}
\usepackage{pgfplots}
\pgfplotsset{compat=1.18}
\usepackage{booktabs,tabularx,array}
\usepackage{enumitem}
\usepackage{fancyhdr}
\usepackage{parskip}
\usepackage{hyperref}
\hypersetup{hidelinks,pdfauthor={Лёвин Артём Александрович},pdftitle={Чек-лист. Задачи прикладного характера. 27 октября 2026}}
\definecolor{accent}{HTML}{008080}
\definecolor{darkaccent}{HTML}{00585A}
\definecolor{ink}{HTML}{333333}
\definecolor{mute}{HTML}{677779}
\color{ink}
\setlength{\parindent}{0pt}
\setlength{\parskip}{5pt}
\renewcommand{\arraystretch}{1.22}
\setlist[itemize]{leftmargin=6mm,itemsep=3pt,topsep=4pt}
\setlist[enumerate]{leftmargin=7mm,itemsep=3pt,topsep=4pt}
\newcommand{\idea}[1]{\textcolor{darkaccent}{\textbf{#1}}}
\newcommand{\result}[1]{\textcolor{darkaccent}{\textbf{Итог:}} #1}
\newcommand{\ex}[2]{\vspace{4pt}\noindent\textcolor{darkaccent}{\textbf{Пример #1. #2}}\par\nopagebreak[3]}
\newcommand{\sectionline}{\par\vspace{-1mm}{\color{accent}\hrule height .45pt}\vspace{4mm}}
\newcommand{\sect}[2]{\section*{\textcolor{darkaccent}{#1. #2}}\sectionline}
\fancypagestyle{study}{%
 \fancyhf{}\fancyhead[L]{\footnotesize Задачи прикладного характера}%
 \fancyhead[R]{\footnotesize Екатерина Гнедкова}%
 \fancyfoot[C]{\small\thepage}%
 \renewcommand{\headrulewidth}{0.45pt}%
 \renewcommand{\headrule}{\hbox to\headwidth{\color{accent}\leaders\hrule height \headrulewidth\hfill}}
 \renewcommand{\footrulewidth}{0pt}}
\pagestyle{study}
\setlength{\headheight}{24pt}
\begin{document}
\thispagestyle{empty}
\begin{center}
\vspace*{24mm}
{\Large\textcolor{darkaccent}{МАТЕМАТИКА}}\par
\vspace{21mm}
{\fontsize{26}{32}\selectfont Чек-лист}\par
\vspace{18mm}
{\fontsize{21}{27}\selectfont\textcolor{accent}{Задачи прикладного\\ характера}}\par
\vspace{15mm}
{\large Подготовка к ЕГЭ по профильной математике}\par
\vfill
{\large Екатерина Гнедкова}\par
\vspace{14mm}
{\large 27 октября 2026 г.}\par
\vspace{18mm}
{\small Лёвин Артём Александрович\\[2mm] эксклюзивно для Екатерины Гнедковой}\par
\vspace{23mm}
\begin{tikzpicture}
\draw[accent,line width=.9pt] (0,0) .. controls (1.4,.55) and (2.2,-.42) .. (3.8,0)
 .. controls (5.3,.48) and (6.7,-.43) .. (8.4,0);
\fill[accent] (3.8,0) circle (1.25pt);
\end{tikzpicture}
\end{center}
\clearpage

\sect{1}{Как устроены прикладные задачи}
\idea{Главная идея.} Текстовое условие описывает процесс, а математическая модель задаёт зависимость между величинами. Нужно понять, что именно сравнивается, составить уравнение или неравенство, найти допустимые решения и перевести их на язык задачи.

\begin{tabularx}{\linewidth}{@{}p{.33\linewidth}X@{}}
\toprule
\textcolor{darkaccent}{Формулировка} & \textcolor{darkaccent}{Математическое действие}\\
\midrule
«увеличилось на», «понизилось на» & Найти разность двух значений величины.\\
«не менее», «не ниже» & Записать неравенство $f(x)\ge a$.\\
«не более», «не выше» & Записать неравенство $f(x)\le a$.\\
«как долго» & Найти длину промежутка: $t_2-t_1$.\\
«наименьшее время» & Выбрать первый допустимый момент достижения границы.\\
«наибольшее расстояние» & Найти правую границу допустимого интервала.\\
\bottomrule
\end{tabularx}

\idea{Проверка перед ответом:}
\begin{enumerate}
\item Указаны ли единицы измерения и приведены ли они к одному виду?
\item Есть ли ограничение $t\ge0$, $x\ge0$ или ограничение времени действия модели?
\item Требуется значение величины, момент времени или продолжительность интервала?
\item Сколько корней получилось и какие из них имеют физический смысл?
\end{enumerate}

\sect{2}{Линейное изменение и разность величин}
\ex{1}{Тепловое расширение рельса}
При $0\,^{\circ}\mathrm{C}$ рельс имеет длину $l_0=10$ м. Его длина при температуре $t$ описывается формулой
\[
 l(t)=l_0(1+\alpha t),\qquad \alpha=1{,}2\cdot10^{-5}\;^{\circ}\mathrm{C}^{-1}.
\]
Длина увеличилась на $3$ мм. Найдём температуру.

Сначала переводим миллиметры в метры: $\Delta l=3\cdot10^{-3}$ м. Из формулы следует
\[
 l-l_0=l_0\alpha t,\qquad
 t=\frac{\Delta l}{l_0\alpha}
 =\frac{3\cdot10^{-3}}{10\cdot1{,}2\cdot10^{-5}}=25.
\]
\result{$25\,^{\circ}\mathrm{C}$. Ключевое действие --- выразить \emph{приращение} $l-l_0$.}

\clearpage
\sect{3}{Изменение уровня воды по времени падения}
\ex{2}{Колодец после дождя}
Расстояние до воды определяется временем падения небольшого камня:
\[
 h=5t^2.
\]
До дождя время составляло $0{,}6$ с, после дождя --- $0{,}4$ с. Сравниваем глубины:
\[
 h_1=5\cdot0{,}6^2=1{,}8\text{ м},\qquad
 h_2=5\cdot0{,}4^2=0{,}8\text{ м}.
\]
Поэтому подъём уровня воды равен
\[
 \Delta h=h_1-h_2=5(t_1^2-t_2^2)=1\text{ м}.
\]
\begin{center}
\begin{tikzpicture}[x=1cm,y=1cm,>=Stealth]
 \draw[darkaccent,line width=1.2pt,rounded corners=5pt] (0,0) -- (0,2.6) -- (1.7,2.6) -- (1.7,0);
 \draw[accent,line width=2pt] (0.06,.75)--(1.64,.75);
 \draw[accent!65,line width=2pt] (0.06,1.45)--(1.64,1.45);
 \draw[<->,thin] (2.15,.75)--(2.15,1.45);
 \node[right,align=left] at (2.3,1.1) {$\Delta h=1$ м};
 \node[left] at (-.12,.75) {до};
 \node[left] at (-.12,1.45) {после};
\end{tikzpicture}
\end{center}
\result{Уровень воды поднялся на $1$ м. Чем меньше время падения, тем ближе поверхность воды.}

\sect{4}{Квадратичная модель: расстояния и высоты}
Если зависимость имеет вид
\[
 f(x)=ax^2+bx+c,\qquad a\ne0,
\]
то график --- парабола. В прикладных задачах часто ищут \emph{интервал}, на котором парабола расположена выше либо ниже заданного уровня.

\ex{3}{Камень должен пролететь над стеной}
Траектория камня описывается формулой
\[
 y=-\frac{x^2}{60}+\frac76x.
\]
Стена имеет высоту $9$ м. Запас высоты должен составлять не менее $1$ м, значит, \idea{нужно $y\ge10$}:
\[
 -\frac{x^2}{60}+\frac76x\ge10
 \quad\Longleftrightarrow\quad x^2-70x+600\le0.
\]
Корни уравнения $x^2-70x+600=0$:
\[
 D=70^2-4\cdot600=2500,\qquad
 x_{1,2}=\frac{70\pm50}{2},\qquad x_1=10,\;x_2=60.
\]
Поскольку квадратный трёхчлен с положительным старшим коэффициентом неположителен \emph{между корнями}, получаем $10\le x\le60$.

\result{Наибольшее расстояние до стены --- $60$ м. Равенство допускается: «не менее» включает границу.}

\clearpage
\sect{5}{Время нахождения выше заданной высоты}
\ex{4}{Подброшенный мяч}
Высота мяча определяется формулой
\[
 h(t)=1{,}4+9t-5t^2,
\]
где $t\ge0$ --- время в секундах. Требуется время, на протяжении которого мяч находится на высоте \emph{не менее} $3$ м.
\[
 1{,}4+9t-5t^2\ge3
 \quad\Longleftrightarrow\quad 5t^2-9t+1{,}6\le0.
\]
Для граничных моментов решаем равенство:
\[
 D=(-9)^2-4\cdot5\cdot1{,}6=49,\qquad
 t_{1,2}=\frac{9\pm7}{10}.
\]
Получаем $t_1=0{,}2$ с и $t_2=1{,}6$ с. Между этими моментами высота не меньше $3$ м.
\[
 \Delta t=t_2-t_1=1{,}6-0{,}2=1{,}4\text{ с}.
\]
\begin{center}
\begin{tikzpicture}
\begin{axis}[width=11cm,height=6.1cm,xmin=0,xmax=1.9,ymin=0,ymax=5.8,
 xlabel={$t$, с},ylabel={$h$, м},xtick={0,0.2,0.9,1.6},
 ytick={0,1.4,3,5.45},grid=major,grid style={black!9},
 axis line style={black!45},tick label style={font=\small},label style={font=\small},
 every axis plot/.append style={thick},clip=false]
 \addplot[domain=0:1.9,samples=160,accent] {1.4+9*x-5*x^2};
 \addplot[domain=0:1.9,samples=2,dashed,black!45] {3};
 \addplot[only marks,mark=*,mark size=2pt,accent] coordinates {(0.2,3) (1.6,3)};
 \draw[{Stealth[length=2mm]}-{Stealth[length=2mm]},darkaccent,line width=.9pt] (axis cs:.2,3.4)--(axis cs:1.6,3.4);
 \node[anchor=south,fill=white,inner sep=2pt,font=\small] at (axis cs:.9,3.4) {$1{,}4$ с};
\end{axis}
\end{tikzpicture}
\end{center}
\result{$1{,}4$ с. Оба корня --- моменты пересечения уровня; длительность равна их разности.}

\clearpage
\sect{6}{Температура: первый момент и продолжительность}
Модель $T(t)=at^2+bt+T_0$ позволяет решать две разные задачи:
\begin{itemize}
\item \idea{когда достигнута граница} --- решаем $T(t)=T_{\text{гр}}$ и выбираем нужный момент;
\item \idea{как долго держится нужная температура} --- решаем неравенство и находим длину промежутка.
\end{itemize}

\ex{5}{Нагреватель с предельной температурой}
\[
 T(t)=1400+200t-10t^2,\qquad T_{\text{гр}}=1760\text{ К}.
\]
Граничные моменты определяются уравнением
\[
 1400+200t-10t^2=1760
 \quad\Longleftrightarrow\quad t^2-20t+36=0,
\]
откуда $t_1=2$ мин и $t_2=18$ мин. На промежутке $(2;18)$ температура \emph{превышает} безопасный предел. Чтобы избежать перегрева, отключаем прибор при первом достижении границы.

\result{Через $2$ минуты после начала работы.}

\ex{6}{Как долго температура не ниже заданной}
\[
 T(t)=-t^2+14t+25,\qquad T(t)\ge70\,^{\circ}\mathrm{C}.
\]
\[
 -t^2+14t+25\ge70
 \quad\Longleftrightarrow\quad t^2-14t+45\le0
 \quad\Longleftrightarrow\quad (t-5)(t-9)\le0.
\]
Значит, $t\in[5;9]$ и искомая продолжительность $9-5=4$ мин.

\result{$4$ минуты. Не путать время достижения границы с длиной интервала.}

\vspace{5mm}
\idea{Разница в постановке вопроса}\par
\begin{tabularx}{\linewidth}{@{}p{.37\linewidth}X@{}}
\toprule
\textcolor{darkaccent}{Требуется} & \textcolor{darkaccent}{Как выбрать результат}\\
\midrule
Первый момент достижения опасной температуры & Меньший допустимый корень уравнения $T(t)=T_{\text{гр}}$.\\
Продолжительность безопасного или рабочего режима & Найти промежуток решения неравенства и вычислить его длину.\\
\bottomrule
\end{tabularx}

\clearpage
\sect{7}{Когда найденные корни имеют физический смысл}
\ex{7}{Уровень воды в баке после открытия крана}
Высота столба воды задаётся формулой
\[
 H(t)=H_0-\sqrt{2gH_0}\,kt+\frac{gk^2t^2}{2},
\]
где $H_0=20$ м, $g=10$ м/с$^2$, $k=\dfrac1{500}$, $t$ --- секунды. Уровень должен снизиться до четверти начального: $H=20/4=5$ м.

Подставляем данные:
\[
 5=20-\frac{20}{500}t+\frac{5}{500^2}t^2
 \quad\Longleftrightarrow\quad t^2-2000t+750000=0.
\]
\[
 D=2000^2-4\cdot750000=1000000,
 \qquad t_1=500,\quad t_2=1500.
\]
Но вода полностью вытечет к моменту, когда $H(t)=0$. В данной модели
\[
 H(t)=5\left(2-\frac{t}{500}\right)^2,
 \qquad H(t)=0\ \Longrightarrow\ t=1000\text{ с}.
\]
После полного опустошения модель уровня столба воды уже не применяется. Поэтому из двух алгебраических корней подходит только $500$ с.

\result{Через $500$ секунд. Допустимый физический промежуток: $0\le t\le1000$ с.}

\ex{8}{Минимальная скорость вращения ведёрка}
Для верхней точки траектории дана формула силы давления воды на дно
\[
 P=m\left(\frac{v^2}{L}-g\right),
\]
где $L=0{,}625$ м, $g=10$ м/с$^2$, а $m>0$ --- масса воды.

При минимальной скорости давление становится равным нулю:
\[
 m\left(\frac{v^2}{L}-g\right)=0.
\]
Поскольку $m>0$, делим на $m$ и получаем
\[
 \frac{v^2}{L}=g,\qquad
 v^2=gL=10\cdot0{,}625=6{,}25,
 \qquad v=\sqrt{6{,}25}=2{,}5\text{ м/с}.
\]
Отрицательный корень скорости исключаем: речь идёт о её \emph{модуле}.

\result{Минимальная скорость $2{,}5$ м/с. Проверка размерности: $\sqrt{(\text{м}/\text{с}^2)\cdot\text{м}}=\text{м}/\text{с}$.}

\clearpage
\begin{center}
{\Large\textcolor{darkaccent}{Алгоритм: от условия к ответу}}\par
\vspace{8mm}
\begin{tikzpicture}[>=Stealth,
 base/.style={draw=darkaccent,line width=.8pt,rounded corners=3pt,align=center,inner sep=7pt,font=\small,fill=white},
 main/.style={base,text width=10.3cm,minimum height=1.25cm},
 branch/.style={base,text width=5.3cm,minimum height=1.5cm},
 question/.style={draw=darkaccent,diamond,aspect=2.5,inner sep=3pt,align=center,font=\small,text width=4.3cm},
 line/.style={-{Stealth[length=2.4mm]},line width=.85pt,darkaccent}]
 \node[main] (start) at (0,0) {Прочитайте условие, выпишите данные\\ и неизвестную величину};
 \node[main] (model) at (0,-2.0) {Переведите единицы, запишите модель\\ $f(x)$ и ограничения на $x$};
 \node[question] (q) at (0,-4.5) {Что требуется найти?};
 \node[branch] (equal) at (-4.15,-7.1) {Момент, расстояние,\\ граничное значение\\\vspace{2mm}$f(x)=a$};
 \node[branch] (ineq) at (4.15,-7.1) {Промежуток, длительность,\\ допустимые значения\\\vspace{2mm}$f(x)\ge a$ или $f(x)\le a$};
 \node[branch] (r1) at (-4.15,-9.7) {Решите уравнение,\\ получите кандидатов};
 \node[branch] (r2) at (4.15,-9.7) {Найдите граничные корни\\ и промежутки нужного знака};
 \node[main] (filter) at (0,-12.25) {Отберите решения по смыслу:\\ $t\ge0$, ограничения модели, единицы измерения};
 \node[main] (end) at (0,-14.35) {Запишите ответ на языке задачи:\\ число, момент, интервал или длительность};
 \draw[line] (start)--(model);
 \draw[line] (model)--(q);
 \draw[line] (q.west)--node[above,sloped,font=\small]{момент/граница} (equal.north);
 \draw[line] (q.east)--node[above,sloped,font=\small]{промежуток} (ineq.north);
 \draw[line] (equal)--(r1);
 \draw[line] (ineq)--(r2);
 \draw[line] (r1.south) -- ++(0,-.48) -| ($(filter.north)+(-2.25,0)$);
 \draw[line] (r2.south) -- ++(0,-.48) -| ($(filter.north)+(2.25,0)$);
 \draw[line] (filter)--(end);
\end{tikzpicture}
\end{center}
\clearpage

\sect{8}{Типичные ошибки и итоговая памятка}
\begin{itemize}
\item \textbf{Единицы.} Три миллиметра равны $3\cdot10^{-3}$ м. Если коэффициент задан в $\text{м}/\text{с}^2$, следите за временем в секундах.
\item \textbf{Разность уровней.} Изменение --- это $h_1-h_2$, а не одно значение $h(t)$; порядок вычитания определяется физическим смыслом.
\item \textbf{Неравенство.} «Не менее» означает знак $\ge$; если речь о времени нахождения выше уровня, сначала ищем границы, затем их разность.
\item \textbf{Второй корень.} Для температуры первое и второе пересечения порога отвечают разным моментам. Для бака поздний корень может выходить за время действия модели.
\item \textbf{Корни и знаки.} При делении неравенства на отрицательное число знак меняется. Для квадратного трёхчлена проверяйте направление ветвей.
\item \textbf{Физические величины.} Масса воды положительна; скорость как модуль неотрицательна. После извлечения корня выбирайте допустимый знак.
\item \textbf{Вычисления.} $10\cdot0{,}625=6{,}25$, $\sqrt{6{,}25}=2{,}5$; оценка порядка величины помогает заметить ошибку.
\end{itemize}

\vspace{4mm}
\begin{tabularx}{\linewidth}{@{}p{.37\linewidth}X@{}}
\toprule
\textcolor{darkaccent}{Сигнал в задаче} & \textcolor{darkaccent}{Первое полезное действие}\\
\midrule
Величина изменилась & Выпишите $\Delta f=f_2-f_1$ с нужным знаком.\\
Дана квадратичная модель & Сравните график с требуемым уровнем.\\
«Не менее» / «не более» & Составьте неравенство.\\
«Сколько времени» & Найдите $t_2-t_1$.\\
«Когда впервые» & Выберите первый допустимый корень.\\
Два решения уравнения & Проверьте область физической применимости.\\
«Минимальная скорость» & Найдите граничный режим и положительный корень.\\
\bottomrule
\end{tabularx}
\vfill
\begin{center}
\small\textcolor{mute}{Лёвин Артём Александрович эксклюзивно для Екатерины Гнедковой}
\end{center}
\end{document}
